\section*{La llegenda del Targus}

\subsection*{I. Clau hermen\`eutica general: la veritat dins del grotesc}
\emph{«En la naturalesa no hi ha res rid\`icul.»} — Blaise Pascal

Aquest ep\`igraf \`es l’\`ancora conceptual del text. El conte proposa una visi\'o del m\'on on fins i tot l’absurd \`es susceptible de significat, si s’interpreta dins una l\`ogica simb\`olica.

\subsection*{II. Tart\`aria: el territori del mite contaminat}
Hist\`oricament vaga i m\'itica, \textbf{Tart\`aria} representa l’espai liminal entre coneixement i caos, mem\`oria i llegenda. \textit{Universitat transformada en regne: lloc d’herois, per\`o tamb\'e de cat\`astrofes.}

\subsection*{III. El Targus: el Desordre Absolut}
\textbf{Targus} \`es entropia viva. Devora semen com a mat\`eria prim\`a del coneixement no estructurat. El seu nom recorda \emph{Tartarus} (infern grec), per\`o tamb\'e \textit{target} — el blanc a abatre.

\textit{Lectures:}
\begin{itemize}
  \item \textbf{Freudiana}: és el desig pur, l’\textit{id} no sublimat.
  \item \textbf{Teol\`ogica}: criatura que castiga l’orgull hum\`a.
  \item \textbf{Cognitiva}: el Drive no classificat. Entropia documental.
\end{itemize}

\subsection*{IV. Seminara: la princesa del l\`imit}
El nom prov\'e de \emph{semen} i \emph{semilla}. \textbf{Seminara} representa el coneixement en estat potencial: prec\`ios, per\`o exposat al perill si no es gestiona.

\subsection*{V. Morvius: el malent\`es tecnocr\`atic}
\textbf{Morvius}, el bruixot, és la figura del \emph{savi incompetent}. Vull fer el bé, però causa el caos. \textit{Al·lusions a l’error tècnic, al científic que toca el que no entén.}

\subsection*{VI. Sant Johan: el gestor del caos}
\textbf{Sant Johan} no \`es nom\'es un heroi: \`es un organitzador. \textit{Amb la seva llan\c{c}a, transforma la viscositat en ordre: neixen les motxilles.}

\subsection*{VII. Les motxilles: objectes sagrats}
S\'on el producte alqu\`imic d’una mort monstruosa. \textit{Simbolitzen la consolidaci\'o funcional del coneixement, l’arxiu purificat.}

\subsection*{VIII. Rei Cl\`amidos: autoritat est\`eril}
El seu nom evoca malaltia ven\`eria. \textbf{Cl\`amidos} \`es el poder que no actua. Figura de la instituci\'o obsoleta, paralitzada davant la crisi.

\subsection*{IX. David Bagan: l’hereu inconscient}
\textbf{David Bagan} representa l’estudiant modern. Hereta la motxilla sense saber-ne el pes simb\`olic. \textit{\`Es l’usuari del coneixement que altres van guanyar.}

\subsection*{X. Cicles i estructura}
\begin{itemize}
  \item \textbf{Exc\`es} $\rightarrow$ \textbf{Caos} $\rightarrow$ \textbf{Combat} $\rightarrow$ \textbf{Sang} $\rightarrow$ \textbf{Objectes} $\rightarrow$ \textbf{Ordre} $\rightarrow$ \textbf{Llegenda}.
\end{itemize}
\textit{El conte \`es una cosmogonia, una g\`enesi institucional disfressada de faula.}

\subsection*{XI. Glossari simb\`olic}
\begin{tabular}{ll}
\textbf{Targus} & Entropia seminal. Caos absolut. \\
\textbf{Seminara} & Coneixement fr\`agil. Passiu. Potencial. \\
\textbf{Morvius} & Error tecnocient\'\i fic. \\
\textbf{Motxilles} & Arxius. Transmissi\'o estructurada. Llegat. \\
\textbf{Sant Johan} & Heroi fundador. Organitzador. \\
\textbf{David Bagan} & Heretador actual. Estudiant. \\
\textbf{Tart\`aria} & Espai del saber ancestral i dif\'\i cil. \\
\end{tabular}

\subsection*{XII. Conclusi\'o}
\textit{La llegenda del Targus} \`es una par\`odia sagrada. Des del grotesc, formula una teoria de la mem\`oria compartida, del coneixement com\`u i del valor d’ordenar l’inabastable.

\vspace{1em}
\textbf{Tota ci\`encia comença en el fang. L’heroi \`es qui el decanta i en fa motxilla.}